% !PRINTED 294
\weberpage{294}
%<<CONT
und sodann, wenn $\begin{pmatrix}\alpha, & \beta\\ \gamma, & \delta\end{pmatrix}$ eine beliebige Substitution in $\mathfrak{L}_0$ und
$\beta$ von $0$ verschieden ist:
\weq{12}{\begin{pmatrix}\alpha, & \beta\\ \gamma, & \delta\end{pmatrix}=\begin{pmatrix}1, & 0\\ \delta\beta^{-1}, & 1\end{pmatrix}\begin{pmatrix}\beta, & 0\\ 0, & \beta^{-1}\end{pmatrix}\begin{pmatrix}0, & 1\\ -1, & 0\end{pmatrix}\begin{pmatrix}1, & 0\\ \alpha\beta^{-1}, & 1\end{pmatrix},}
und wenn $\beta=0$ ist:
\weq{13}{\begin{pmatrix}\alpha, & 0\\ \gamma, & \alpha^{-1}\end{pmatrix}=\begin{pmatrix}\alpha, & 0\\ 0, & \alpha^{-1}\end{pmatrix}\begin{pmatrix}1, & 0\\ \gamma\alpha, & 1\end{pmatrix}.}
\gesperrt{Es ist aus diesem Satz zu schließen, daß eine in $\mathfrak{L}_0$ enthaltene
 Gruppe, die die zwei Substitutionen $A, B$ enthält,
notwendig mit $\mathfrak{L}_0$ identisch sein muß.}

\begin{center}
\S\,80. \gesperrt{Normalteiler der Gruppe $\mathfrak{L}_0$.}
\end{center}

Um die etwa möglichen Reduktionen der Transformationsgleichung
 kennen zu lernen, ist vor allem erforderlich, die Divisoren
 der Gruppe $\mathfrak{L}_0$ zu untersuchen. Wir fragen zuerst nach
der Existenz eines \gesperrt{Normalteilers} $\mathfrak{R}$ von $\mathfrak{L}_0$ (Bd. II, \S\,3)\footnote{Nach \gesperrt{Galois}: \glqq Eigentliche Teiler\grqq{}.}.

Ist
\weq{1}{S=\begin{pmatrix}\alpha, & \beta\\ \gamma, & \delta\end{pmatrix}}
irgend eine von der identischen Substitution verschiedene Substitution
 in $\mathfrak{R}$, so ist, nach dem Wesen des Normalteilers, jede
Substitution
\weq{2}{U=TST^{-1}}
gleichfalls in $\mathfrak{R}$ enthalten, wenn
\weq{3}{T=\begin{pmatrix}\alpha', & \beta'\\ \gamma', & \delta'\end{pmatrix}}
eine \gesperrt{beliebige} Substitution in $\mathfrak{L}_0$ ist.
Wir stellen uns die Aufgabe, $T$ und $\xi$ so zu bestimmen, daß
\weq{4}{U=\begin{pmatrix}0, & 1\\ -1, & \xi\end{pmatrix}}
wird. Diese Bedingung kann auch so geschrieben werden:
\[
\begin{pmatrix}\alpha', & \beta'\\ \gamma', & \delta'\end{pmatrix}\begin{pmatrix}\alpha, & \beta\\ \gamma, & \delta\end{pmatrix}=\begin{pmatrix}0, & 1\\ -1, & \xi\end{pmatrix}\begin{pmatrix}\alpha', & \beta'\\ \gamma', & \delta'\end{pmatrix},
\]
%CONT>>
